\chapter{STUB Mode Target Flow}

\section*{Purpose}
This chapter defines the first real target-hardware bring-up flow, explicitly limited to `STUB` mode.

\section{Why STUB Mode Is First}
The deterministic STUB engine mode removes ML-DSA datapath risk from the first board interaction. This allows bring-up to focus on the PS-visible AXI-Lite wrapper contract, Linux MMIO access, address assignment, and command sequencing.

\section{Required Top-Level Selection}
The first board image shall use `pqsig\_top\_stub\_mode.sv`, which instantiates the stable wrapper with `ENGINE\_MODE\_STUB`. There is no software-visible runtime mode switch in the current design.

\section{Target Flow}
\begin{enumerate}[leftmargin=*]
  \item Build a bitstream using the STUB-mode top-level scaffold.
  \item Load the bitstream on the target Zynq platform.
  \item Boot Linux on the PS.
  \item Provide the AXI-Lite wrapper base address to the software real backend.
  \item Run the safe MMIO probe path.
  \item Run the explicit STUB selftest path.
  \item Confirm the returned signature matches the documented deterministic STUB rule.
\end{enumerate}

\section{Expected Software Commands}
Example software flow on the target board:
\begin{itemize}[leftmargin=*]
  \item `python -m sw.daemon.main probe-mmio --backend real --mmio-base-addr <addr>`
  \item `python -m sw.daemon.main selftest --backend real --mmio-base-addr <addr>`
\end{itemize}

\section{Success Criteria}
Bring-up success for this phase means the wrapper is visible from Linux and one STUB-mode selftest transaction completes with the exact deterministic fake-signature result.

\section{Deferred Transition}
Only after STUB-mode board validation is complete should the project move to an MLDSA\_OSH-mode target image. That later phase must preserve the wrapper contract and software MMIO interface while expanding the hardware internals behind the adapter boundary.